\section{问题分析}
四个问题共享同一套微网物理结构：小区负荷必须被满足，光伏优先用于当前负荷，储能承担跨时段能量转移，外部电网负责补充正常购电或紧急购电。四问之间真正变化的是\textbf{决策时刻能够获得的信息量}以及\textbf{购电结算规则}。因此不宜将四问分别建立完全独立的模型，而应从问题一的显式能量分流模型出发逐层扩展。

\subsection{问题一分析}

问题一中，当天电价、负荷和光伏预测曲线均视为已知，属于典型的单日确定性优化问题。一天内电价存在明显峰谷差，而光伏出力集中在白天，且部分时段存在光伏超过负荷的情形。因此储能具有两种直接作用：一是吸收富余光伏，二是在低价时段利用外网充电、在高价时段放电。

为同时描述外网供负荷、外网充电、光伏充电、储能放电和弃光的去向，需要将总能量平衡进一步拆分为显式流量平衡。又因为同一时段储能不能同时充电和放电，需要引入 0-1 状态变量，因此问题一适合建立\textbf{能量分流混合整数线性规划模型（MILP）}，以全天购电费用最小为目标。

问题一的核心作用不仅是求出一个该日策略，更重要的是形成后续三问可以直接复用的微网物理层。

\subsection{问题二分析}

问题二相对于问题一的根本变化是：每天 0:00 制定购电计划时，当天真实负荷和真实光伏并未全部实现，而计划购电不足又会触发 5 倍电价的紧急购电。

因此，问题二不能简单地把问题一模型按 334 天重复求解，也不能在 0:00 把储能充放电等实际执行变量全部锁死。真正需要提前确定的是\textbf{计划购电量}，而储能充放电、弃光、紧急购电等应当在随机情景实现后再确定。

由此自然形成两层决策：

\[
\boxed{
\text{第一阶段：0:00 确定计划购电}
\quad\longrightarrow\quad
\text{第二阶段：情景实现后进行物理补救}
}
\]

同时，为避免单日模型在当天末过度放电，需要使当前日末储电量具有未来价值。因此采用 7 日滚动视野，并利用历史预测残差构造 SAA 情景，将预测误差风险直接纳入日前计划。问题二由此建立为\textbf{7 日滚动 SAA 两阶段随机 MILP}。

\subsection{问题三分析}

问题三中微网设备和购电物理过程没有变化，新增的是 6:00、12:00、18:00 三次信息更新机会。每天 0:00 已有原始计划，但新光伏预报到达后，尚未执行的计划可以重新调整。

因此问题三的关键不是重新设计能源流模型，而是解决三个问题：

\begin{enumerate}
\item 新预报如何替换旧预报；
\item 已经执行的历史时段如何锁定；
\item 调增、调减与紧急购电如何统一结算。
\end{enumerate}

题目给出的光伏预报仍然只是预测而非真值，因此最新预报只能作为中心预测，仍需叠加历史预报误差构造随机情景。每次更新时，应使用当前真实储电量作为新阶段初值，并仅对尚未执行区间重新优化。

因此问题三是在问题二模型基础上引入阶段集合

\[
\mathcal S=\{0,6,12,18\},
\]

形成\textbf{基于多时点光伏预报更新的滚动 SAA-MILP 购电调整模型}。

\subsection{问题四分析}

问题四进一步取消“未来电价已知”的假设。此时负荷、光伏和电价均存在不确定性。未来真实电价只能在其实际发生后用于结算，不能提前进入计划优化，否则会造成未来信息泄漏。

对于 Q4-2，需要在问题二模型中增加未来电价预测，并将价格预测误差与负荷、光伏误差尽量按同一历史日联合抽样，形成

\[
(L^{(\omega)},PV^{(\omega)},\pi^{(\omega)})
\]

联合情景。

对于 Q4-3，还需要在问题三的 0:00、6:00、12:00、18:00 四阶段滚动结构中，根据已经实现的真实价格对剩余时段的价格预测进行因果更新。

此外，历史价格误差叠加后可能产生负价格情景。若仍允许无上界的计划购电，模型可能出现“为了负成本而购买无实际用途电量”的数学套利；若只用连续变量表示调增和调减，负价格还可能导致二者同时为正。因此问题四还需要额外加入：

\[
\text{H-1：计划购电物理上界}
\]

和

\[
\text{H-2：调增/调减显式互斥}.
\]

这两项约束是问题四相对于问题二、三最重要的针对性扩展。

\subsection{四问统一建模关系}

四问并非四个彼此独立的模型，而是同一物理系统上的逐层递进：

\[
\boxed{
\begin{array}{c}
\text{Q1：确定性能量分流 MILP}\\
\downarrow\\
\text{Q2：计划变量与执行变量分层 + SAA + 7 日滚动}\\
\downarrow\\
\text{Q3：0/6/12/18 信息更新 + 计划再调整}\\
\downarrow\\
\text{Q4：价格预测 + 联合价格情景 + 负价加固}
\end{array}
}
\]

因此全文统一复用问题一建立的物理流量变量，只在后续问题中根据新增业务规则增加必要变量。这样可以保证模型之间的连续性、符号的一致性和结果的可解释性。

% 待补：建模流程图
% 初稿在摘要后留有「建模流程图」占位（四问由确定性调度逐步扩展至联合不确定性优化）。
% 该图属建模手交付物，不在 figures/ 内，故此处只登记、不插入 figure，避免编译时找不到图片。
